\chapter{Stochastic differential equations}
\label{ch:sde}

\begin{question}
If the drift and the noise of a process depend on the process itself, can we solve for it? Which interest-rate and price models come out?
\end{question}

\begin{whatfor}
\begin{itemize}
  \item Defines a stochastic differential equation (SDE) and its solution.
  \item Solves the Ornstein--Uhlenbeck SDE with an integrating factor and finds its normal distribution.
  \item Applies it to the Vasicek interest-rate model; states the CIR model.
  \item Solves the geometric Brownian motion SDE $dN=rN\,dt+\alpha N\,dW$.
\end{itemize}
\skipline{fs:sde}
\end{whatfor}

\begin{prereq}
From Part~0: normal distribution, $\Phi$. From \crefrange{ch:gbm}{ch:ito-formula}: geometric Brownian motion, Gaussian process, Itô integral of a non-random function is $N(0,\int f^2)$, Itô isometry, mean $0$ of Itô integrals, Itô process, Itô's formula, product rule.
\end{prereq}

%=====================================================================
\section{SDEs and their solutions}

\begin{definition}[Stochastic differential equation]
\label{def:sde}
Let $a(t,x)$ and $b(t,x)$ be non-random functions. A process $X$ \dfn{solves the SDE}
\[
  dX_t=a(t,X_t)\,dt+b(t,X_t)\,dW_t,\qquad X_0=x_0,
\]
if it is an Itô process with $X_t=x_0+\int_0^ta(s,X_s)\,ds+\int_0^tb(s,X_s)\,dW_s$ for all $t\in[0,T]$.
\end{definition}

The difference from an Itô process: the coefficients are built from the unknown $X$ itself, so this is an equation to solve. For example $dX_t=dt+2\,dW_t$, $X_0=3$ has the solution $X_t=3+t+2W_t$ (just integrate), while $dX_t=-2X_t\,dt+dW_t$ cannot be integrated directly: $X$ appears on both sides.

SDEs rarely have a closed-form solution. Two that do are the main models of this chapter. The tool for both is Itô's formula with a well-chosen $f$, chosen so that the unknown disappears.

%=====================================================================
\section{The Ornstein--Uhlenbeck process}

\begin{definition}[Ornstein--Uhlenbeck process]
For constants $r,\sigma>0$, the solution of
\[
  dX_t=-rX_t\,dt+\sigma\,dW_t,\qquad X_0=x_0,
\]
is the \dfn{Ornstein--Uhlenbeck process} (OU). The drift $-rX_t$ always points back towards $0$, and is stronger the farther $X$ is from $0$: \dfn{mean reversion}.
\end{definition}

\begin{theorem}[Solution of the OU SDE]
\label{thm:ou}
\[
  X_t=x_0e^{-rt}+\sigma\int_0^te^{-r(t-s)}\,dW_s,
  \qquad
  X_t\sim N\Big(x_0e^{-rt},\ \frac{\sigma^2}{2r}\big(1-e^{-2rt}\big)\Big).
\]
\end{theorem}
\begin{proof}
Multiply by the integrating factor $e^{rt}$.
\begin{steps}
  \item Itô with $f(t,x)=e^{rt}x$ ($\partial_tf=re^{rt}x$, $\partial_xf=e^{rt}$, $\partial_{xx}f=0$):
  \[
    d\big(e^{rt}X_t\big)=re^{rt}X_t\,dt+e^{rt}dX_t=re^{rt}X_t\,dt+e^{rt}\big(-rX_t\,dt+\sigma\,dW_t\big)=\sigma e^{rt}\,dW_t .
  \]
  The unknown $X$ has cancelled.
  \item Integrate: $e^{rt}X_t-x_0=\sigma\int_0^te^{rs}\,dW_s$. Multiply by $e^{-rt}$ for the solution.
  \item The integrand $e^{-r(t-s)}$ is non-random, so $Q_t=\int_0^te^{-r(t-s)}dW_s$ is normal with mean $0$ and, by the isometry,
  \[
    \Var Q_t=\int_0^te^{-2r(t-s)}\,ds=e^{-2rt}\Big[\frac{e^{2rs}}{2r}\Big]_0^t=\frac{1-e^{-2rt}}{2r}.
  \]
  \item $X_t=x_0e^{-rt}+\sigma Q_t$: mean $x_0e^{-rt}$, variance $\sigma^2\Var Q_t$. \qedhere
\end{steps}
\end{proof}

As $t\to\infty$, $X_t$ tends in distribution to $N\big(0,\frac{\sigma^2}{2r}\big)$: the start is forgotten and the spread stays bounded, unlike Brownian motion whose variance grows like $t$. If $X_0$ is normal and independent of $W$, then $X$ is a Gaussian process.

\begin{pic}
\begin{tikzpicture}
\begin{axis}[width=10.5cm,height=5cm,xmin=0,xmax=1,xlabel={$t$},ylabel={$X_t$},samples=100,domain=0:1]
  \addplot[draw=none,fill=thmc!12] {2*exp(-2*x)+sqrt(0.25*(1-exp(-4*x)))} \closedcycle;
  \addplot[draw=none,fill=white] {2*exp(-2*x)-sqrt(0.25*(1-exp(-4*x)))} \closedcycle;
  \addplot[defc,thick] table[x=t,y=x1] {figs/ou.dat};
  \addplot[trapc,thick] table[x=t,y=x2] {figs/ou.dat};
  \addplot[prec,thick] table[x=t,y=x3] {figs/ou.dat};
  \addplot[black,very thick,dashed] {2*exp(-2*x)};
  \addplot[prec,dotted] {0};
\end{axis}
\end{tikzpicture}
\picnote{Three OU paths with $x_0=2$, $r=2$, $\sigma=1$. They are pulled back towards $0$ along the mean $2e^{-2t}$ (dashed); the band is one standard deviation, which levels off at $\sigma/\sqrt{2r}=0.5$.}
\end{pic}

%=====================================================================
\section{Interest-rate models}

\begin{definition}[Vasicek and CIR models]
For constants $a,b,\sigma>0$, the short interest rate $r_t$ follows the
\begin{itemize}
  \item \dfn{Vasicek model} if $dr_t=a(b-r_t)\,dt+\sigma\,dW_t$;
  \item \dfn{CIR model} (Cox--Ingersoll--Ross) if $dr_t=a(b-r_t)\,dt+\sigma\sqrt{r_t}\,dW_t$, $r_0>0$.
\end{itemize}
In both, $b$ is the long-run level and $a$ the speed of return to it.
\end{definition}

\begin{result}[The Vasicek rate is normal]
\label{res:vasicek}
$r_t\sim N\Big(b+e^{-at}(r_0-b),\ \dfrac{\sigma^2}{2a}\big(1-e^{-2at}\big)\Big)$.
\end{result}
\begin{proof}
Put $X_t=r_t-b$. Then $dX_t=dr_t=-aX_t\,dt+\sigma\,dW_t$: an OU process with $r\to a$ and $x_0=r_0-b$. Apply \cref{thm:ou} and add $b$.
\end{proof}

\begin{example}[Can the rate go negative?]
\label{ex:vasicek}
$a=0.5$, $b=0.05$, $\sigma=0.02$, $r_0=0.03$. At $t=2$: mean $0.05+e^{-1}(-0.02)\approx0.0426$, variance $\frac{0.0004}{1}(1-e^{-2})\approx0.000346$, s.d.\ $\approx0.0186$. So $P(r_2<0)=\Phi(-0.0426/0.0186)\approx\Phi(-2.29)\approx0.011$. A normal rate can be negative; that is the weakness the CIR model fixes.
\end{example}

In the CIR model the noise $\sigma\sqrt{r_t}$ dies out as the rate approaches $0$, while the drift $ab>0$ pushes it up, so $r_t$ stays non-negative. The integrating factor still gives $r_t=b+e^{-at}(r_0-b)+\sigma\int_0^te^{-a(t-s)}\sqrt{r_s}\,dW_s$, which is not a closed form (the unknown is inside the integral), but the mean is the same as Vasicek's: $\E_P(r_t)=b+e^{-at}(r_0-b)$, since the Itô integral has mean $0$.

%=====================================================================
\section{Geometric Brownian motion as an SDE}

\begin{theorem}[Solution of the GBM SDE]
\label{thm:gbm-sde}
For constants $r$ and $\alpha>0$, the SDE $dN_t=rN_t\,dt+\alpha N_t\,dW_t$, $N_0>0$, has the solution
\[
  N_t=N_0\exp\Big\{\Big(r-\frac{\alpha^2}2\Big)t+\alpha W_t\Big\},\qquad \E_P(N_t)=N_0e^{rt}.
\]
\end{theorem}
\begin{proof}
Take logarithms with Itô.
\begin{steps}
  \item $f(x)=\ln x$: $f'=\frac1x$, $f''=-\frac1{x^2}$, and $(dN_t)^2=\alpha^2N_t^2\,dt$. So
  \[
    d(\ln N_t)=\frac{dN_t}{N_t}-\frac{1}{2N_t^2}\alpha^2N_t^2\,dt=\Big(r-\frac{\alpha^2}2\Big)dt+\alpha\,dW_t .
  \]
  \item Integrate: $\ln N_t-\ln N_0=\big(r-\frac{\alpha^2}{2}\big)t+\alpha W_t$; exponentiate.
  \item Check (this also covers the assumption $N_t>0$ made by taking $\ln$): by the GBM example of \cref{ch:ito-formula}, $e^{\mu t+\sigma W_t}$ has $dt$-coefficient $(\mu+\frac{\sigma^2}{2})$ times itself; here $\mu+\frac{\sigma^2}2=r$.
  \item Mean: $\E_PN_t=N_0e^{(r-\alpha^2/2)t}\E_Pe^{\alpha W_t}=N_0e^{(r-\alpha^2/2)t}e^{\alpha^2t/2}=N_0e^{rt}$. \qedhere
\end{steps}
\end{proof}

The equation $\frac{dN_t}{N_t}=r\,dt+\alpha\,dW_t$ is just notation for the SDE: the percentage change has mean rate $r$ and volatility $\alpha$. The median of $N_t$ is $N_0e^{(r-\alpha^2/2)t}$, below the mean $N_0e^{rt}$.

%=====================================================================
\section{A volatility model}

\begin{example}[The Heston variance process]
\label{ex:heston}
A ``variance process'' follows $dV_t=(1-V_t)\,dt+2\sqrt{V_t}\,dW_t$, $V_0=1$. (i) Find the SDE of the volatility $Z_t=\sqrt{V_t}$ and $Z_0$. (ii) Show that $Z_t=e^{-t/2}\big(1+\int_0^te^{s/2}dW_s\big)$ solves it.
\begin{steps}
  \item (i) $f(x)=\sqrt x$: $f'=\frac12x^{-1/2}$, $f''=-\frac14x^{-3/2}$; $(dV_t)^2=4V_t\,dt$. Itô:
  \[
    dZ_t=\frac{(1-V_t)\,dt+2\sqrt{V_t}\,dW_t}{2\sqrt{V_t}}-\frac{4V_t\,dt}{8V_t^{3/2}}=dW_t-\frac12\sqrt{V_t}\,dt=dW_t-\frac12Z_t\,dt,
  \]
  with $Z_0=\sqrt{V_0}=1$. (The terms $\pm\frac{1}{2\sqrt{V_t}}dt$ cancel.)
  \item (ii) $Z_t=X_tY_t$ with $X_t=e^{-t/2}$ ($dX_t=-\frac12X_t\,dt$, non-random) and $Y_t=1+\int_0^te^{s/2}dW_s$ ($dY_t=e^{t/2}dW_t$). Product rule, with $dX_t\,dY_t=0$:
  \[
    dZ_t=X_t\,dY_t+Y_t\,dX_t=e^{-t/2}e^{t/2}\,dW_t-\frac12X_tY_t\,dt=dW_t-\frac12Z_t\,dt,
  \]
  and $Z_0=1\cdot1=1$. It is an OU process with $r=\frac12$, $\sigma=1$.
\end{steps}
\end{example}

%=====================================================================
\formulasheet{fs:sde}

Geometric Brownian motion $dN=\mu N\,dt+\sigma N\,dW$ is Card~\ref{card:gbm} with $m=\mu-\frac12\sigma^2$. This sheet holds the mean-reverting models.

\begin{mastercard}[label=card:linsde]{The linear (mean-reverting) SDE}
For constants $\alpha$, $\beta>0$, $\sigma>0$, the SDE $dX_t=(\alpha-\beta X_t)\,dt+\sigma\,dW_t$, $X_0=x$, is solved by the integrating factor $e^{\beta t}$ (so $d(e^{\beta t}X_t)=\alpha e^{\beta t}dt+\sigma e^{\beta t}dW_t$):
\[
  \boxed{\begin{gathered}X_t=xe^{-\beta t}+\frac{\alpha}{\beta}\big(1-e^{-\beta t}\big)+\sigma\int_0^te^{-\beta(t-s)}\,dW_s,\\
  X_t\sim N\Big(xe^{-\beta t}+\frac\alpha\beta\big(1-e^{-\beta t}\big),\ \frac{\sigma^2}{2\beta}\big(1-e^{-2\beta t}\big)\Big).\end{gathered}}
\]
The mean moves from $x$ towards the long-run level $\alpha/\beta$; the variance rises to $\sigma^2/(2\beta)$ and stays bounded.
\begin{plugrows}
  \plugrow{Ornstein--Uhlenbeck}{\pstep{Box: $X_t\sim N\Big(xe^{-\beta t}+\frac\alpha\beta(1-e^{-\beta t}),\ \frac{\sigma^2}{2\beta}(1-e^{-2\beta t})\Big)$.}\pstep{From $dX=-rX\,dt+\sigma\,dW$: $\alpha=0$, $\beta=r$.}\pstep{Substitute; for $t\to\infty$, $e^{-rt}\to0$.}}{$X_t\sim N\big(xe^{-rt},\ \frac{\sigma^2}{2r}(1-e^{-2rt})\big)$; $t\to\infty$: $N(0,\frac{\sigma^2}{2r})$}{$r>0$}{$dX=-X\,dt+dW$ ($r=1$, $\sigma=1$): $\Var X_t=\frac12(1-e^{-2t})$; $x=2$, $r=2$, $\sigma=1$, $t=1$: mean $2e^{-2}=0.271$, variance $\frac14(1-e^{-4})=0.245$}
  \plugrow{Vasicek short rate}{\pstep{Box: $X_t\sim N\Big(xe^{-\beta t}+\frac\alpha\beta(1-e^{-\beta t}),\ \frac{\sigma^2}{2\beta}(1-e^{-2\beta t})\Big)$ for $dX=(\alpha-\beta X)dt+\sigma dW$.}\pstep{From $dr=a(b-r)dt+\sigma dW=(ab-ar)dt+\sigma dW$: $\alpha=ab$, $\beta=a$, $x=r_0$.}\pstep{Substitute: mean $r_0e^{-at}+b(1-e^{-at})=b+e^{-at}(r_0-b)$.}}{$r_t\sim N\big(b+e^{-at}(r_0-b),\ \frac{\sigma^2}{2a}(1-e^{-2at})\big)$}{rates can go negative}{$a=0.5$, $b=0.05$, $\sigma=0.02$, $r_0=0.03$, $t=2$: mean $0.05-0.02e^{-1}=0.0426$, variance $\frac{0.0004}{1}(1-e^{-2})=0.000346$, s.d.\ $0.0186$; $P(r_2<0)=\Phi\big(\frac{-0.0426}{0.0186}\big)=\Phi(-2.29)\approx0.011$}
  \plugrow{CIR short rate}{\pstep{$dr=a(b-r)dt+\sigma\sqrt r\,dW$: the noise is not constant, so the box's normal law does not apply.}\pstep{Take means: the $dW$ part has mean $0$ (\cref{card:itoint}), so $\frac d{dt}\E_Pr_t=a(b-\E_Pr_t)$.}\pstep{The same equation as Vasicek's mean, so the same solution.}}{$\E_Pr_t=b+e^{-at}(r_0-b)$ (same as Vasicek); $r_t\ge0$, not normal}{no closed-form law}{$a=0.5$, $b=0.05$, $r_0=0.03$: $\E_Pr_2=0.05-0.02e^{-1}\approx0.0426$}
  \plugrow{A volatility process}{\pstep{Box, first line: $X_t=xe^{-\beta t}+\frac\alpha\beta(1-e^{-\beta t})+\sigma\int_0^te^{-\beta(t-s)}dW_s$.}\pstep{From $dZ=-\frac12Z\,dt+dW$, $Z_0=1$: $\alpha=0$, $\beta=\frac12$, $\sigma=1$, $x=1$.}\pstep{Substitute and take $e^{-t/2}$ out.}}{$Z_t=e^{-t/2}\big(1+\int_0^te^{s/2}dW_s\big)$}{came from $Z=\sqrt V$ by Card~\ref{card:ito}}{variance from the box: $\frac{1}{2\cdot\frac12}(1-e^{-t})=1-e^{-t}$}
\end{plugrows}
\end{mastercard}
